% prompt2analytics Lehrfolien (Deutsch/Schweiz)
% Beamer mit Metropolis-Theme
\documentclass[aspectratio=169]{beamer}
\usetheme{metropolis}

% Pakete
\usepackage[T1]{fontenc}
\usepackage[utf8]{inputenc}
\usepackage[ngerman]{babel}  % Schweizer Deutsch
\usepackage{booktabs}
\usepackage{graphicx}
\usepackage{listings}
\usepackage{hyperref}

% Listings-Konfiguration für Code-Beispiele
\lstset{
    basicstyle=\ttfamily\small,
    breaklines=true,
    frame=single,
    backgroundcolor=\color{gray!10}
}

% Titelinformationen
\title{Einführung in prompt2analytics}
\subtitle{Datenanalyse durch natürliche Sprache}
\author{[Dozent/in]}
\institute{[Institution]}
\date{\today}

\begin{document}

% Titelfolie
\maketitle

% Inhaltsverzeichnis
\begin{frame}{Übersicht}
    \tableofcontents
\end{frame}

% =============================================================================
\section{Einführung}
% =============================================================================

\begin{frame}{Was ist prompt2analytics?}
    % TODO: Inhalt ergänzen
    \begin{itemize}
        \item Ein modernes Analyse-Toolkit
        \item Natürlichsprachliche Schnittstelle für statistische Analysen
        \item Reine Rust-Implementierung für hohe Performance
    \end{itemize}
\end{frame}

\begin{frame}{Warum p2a verwenden?}
    % TODO: Inhalt ergänzen
    \begin{itemize}
        \item Keine Programmierkenntnisse für Chat-Interface nötig
        \item Reproduzierbare Analyse-Workflows
        \item Professionelle Ökonometrie- und ML-Funktionen
    \end{itemize}
\end{frame}

% =============================================================================
\section{Erste Schritte}
% =============================================================================

\begin{frame}{Installation}
    % TODO: Installationsanleitung ergänzen
\end{frame}

\begin{frame}{Ersten Datensatz laden}
    % TODO: Inhalt mit Screenshots ergänzen
\end{frame}

% =============================================================================
\section{Das Chat-Interface}
% =============================================================================

\begin{frame}{Übersicht der Web-Oberfläche}
    % TODO: Screenshot und Beschreibung ergänzen
\end{frame}

\begin{frame}{Natürlichsprachliche Abfragen}
    % TODO: Beispiele ergänzen
    \begin{itemize}
        \item «Lade die Verkaufsdaten»
        \item «Zeig mir eine Zusammenfassung des Datensatzes»
        \item «Führe eine Regression von Umsatz auf Werbeausgaben durch»
    \end{itemize}
\end{frame}

\begin{frame}{Ergebnisse anzeigen}
    % TODO: Inhalt zum Ergebnisbereich ergänzen
\end{frame}

% =============================================================================
\section{CLI-Grundlagen}
% =============================================================================

\begin{frame}[fragile]{Befehlsstruktur}
    % TODO: CLI-Beispiele ergänzen
    \begin{lstlisting}
p2a data load meinedaten.csv --name verkauf
p2a data describe verkauf
p2a reg ols verkauf -y umsatz -x werbung
    \end{lstlisting}
\end{frame}

% =============================================================================
\section{Praktische Beispiele}
% =============================================================================

\begin{frame}{Beispiel: Verkaufsanalyse}
    % TODO: Durchgang ergänzen
\end{frame}

\begin{frame}{Beispiel: Wirtschaftsdatenanalyse}
    % TODO: Durchgang ergänzen
\end{frame}

% =============================================================================
\section{Zusammenfassung}
% =============================================================================

\begin{frame}{Wichtigste Punkte}
    % TODO: Zusammenfassung ergänzen
\end{frame}

\begin{frame}{Ressourcen}
    \begin{itemize}
        \item Dokumentation: \url{https://github.com/yourusername/prompt2analytics}
        \item Tutorials: Siehe Ordner \texttt{tutorials/de-CH/}
        \item Beispieldaten: Siehe Ordner \texttt{teaching\_data/}
    \end{itemize}
\end{frame}

\end{document}
